\section{Recensement des signatures de l’émetteur résiduel}
\label{app:firmas-emision}

Ce recensement spécifie les fibres finies utilisées dans la section~\ref{sec:extension}.
Chaque condition initiale est identifiée par
\[
 \nu=4(9i+j)+d,\qquad 0\leq i,j\leq8,\quad
 d=0,1,2,3\ \longleftrightarrow\ N,E,S,O.
\]
La division euclidienne retrouve \(d=[\nu]_4\), \(j=[\lfloor\nu/4\rfloor]_9\)
et \(i=\lfloor\nu/36\rfloor\). Par conséquent, \(0\leq\nu<324\) énumère
chaque condition exactement une fois, et les paires de ces identifiants
énumèrent les \(104\,976\) conditions de \(\Omega_{\rm seed}\).

Soit \(f_+(\nu)\), respectivement \(f_\times(\nu)\), la signature à six
composantes obtenue par la récurrence d’émission. À chaque étape active, on lit
la case avant d’avancer ; les vecteurs de direction sont
\[
 v_N=(-1,0),\quad v_E=(0,1),\quad v_S=(1,0),\quad v_O=(0,-1).
\]
Les coordonnées sont réduites modulo neuf et le vecteur est inversé après les étapes \(27\) et \(54\). Les étapes de phase zéro ne déplacent aucun des curseurs. Ces règles, avec \eqref{eq:emisor-seis}, déterminent \(f_+\) et \(f_\times\) sans catalogue de valeurs cibles.

Le tableau donne toutes leurs images. La multiplicité \(m_\sigma\) est le nombre
d’identifiants dans la fibre \(f_\sigma^{-1}(a)\) ; \(\nu_{\min}\) est son
plus petit élément. Chaque mot conserve les zéros initiaux. La fibre complète
est déterminée par la condition
\[
 f_\sigma^{-1}(a)=\{\nu\in\{0,\ldots,323\}:f_\sigma(\nu)=a\}.
\]
\begin{center}\small
\begin{tabular}{@{}lrr@{\qquad}lrr@{}}
\toprule
\(f_+\) & \(m_+\) & \(\nu_{\min}\) & \(f_\times\) & \(m_\times\) & \(\nu_{\min}\)\\
\midrule
\texttt{122564} & 18 & 17 & \texttt{000000} & 108 & 8 \\
\texttt{122898} & 18 & 24 & \texttt{177393} & 8 & 1 \\
\texttt{212645} & 18 & 5 & \texttt{177555} & 8 & 54 \\
\texttt{212988} & 18 & 0 & \texttt{177771} & 8 & 11 \\
\texttt{221456} & 18 & 29 & \texttt{339555} & 8 & 6 \\
\texttt{221889} & 18 & 12 & \texttt{339771} & 8 & 15 \\
\texttt{456221} & 18 & 28 & \texttt{339933} & 8 & 59 \\
\texttt{465898} & 18 & 21 & \texttt{393177} & 8 & 0 \\
\texttt{546988} & 18 & 9 & \texttt{393393} & 8 & 45 \\
\texttt{564122} & 18 & 16 & \texttt{393555} & 8 & 40 \\
\texttt{645212} & 18 & 4 & \texttt{555177} & 8 & 52 \\
\texttt{654889} & 18 & 33 & \texttt{555339} & 8 & 4 \\
\texttt{889221} & 18 & 13 & \texttt{555393} & 8 & 41 \\
\texttt{889654} & 18 & 32 & \texttt{555555} & 24 & 84 \\
\texttt{898122} & 18 & 25 & \texttt{555717} & 8 & 14 \\
\texttt{898465} & 18 & 20 & \texttt{555771} & 8 & 18 \\
\texttt{988212} & 18 & 1 & \texttt{555933} & 8 & 24 \\
\texttt{988546} & 18 & 8 & \texttt{717555} & 8 & 12 \\
 & &  & \texttt{717717} & 8 & 21 \\
 & &  & \texttt{717933} & 8 & 19 \\
 & &  & \texttt{771177} & 8 & 9 \\
 & &  & \texttt{771339} & 8 & 13 \\
 & &  & \texttt{771555} & 8 & 16 \\
 & &  & \texttt{933339} & 8 & 57 \\
 & &  & \texttt{933555} & 8 & 26 \\
 & &  & \texttt{933717} & 8 & 17 \\
\bottomrule
\end{tabular}
\end{center}

L’évaluation de la récurrence pour les \(324\) identifiants produit
les deux partitions du tableau : \(18\cdot18=324\) et
\(24\cdot8+24+108=324\). Le produit de leurs fibres donne
\(432\) fibres de taille \(144\), \(18\) de taille \(432\) et \(18\)
de taille \(1944\), par l’injectivité de \eqref{eq:acoplamiento}.
Enfin, appliquer \eqref{eq:psi-catalogo} à ce produit d’images
donne \(243\) mots distincts. Le recensement est fini et appartient à l’émetteur
résiduel ; le raffinement coinductif ultérieur conserve sa propre loi.

